\begin{figure}[tbp]
\centering
\begin{tikzpicture}
\begin{loglogaxis}[tlplot, width=0.62\linewidth, height=0.42\linewidth,
  xlabel={lag $L$ (m)}, ylabel={$\nu(L)$ (m)},
  xtick={270,360,540,720,1080}, xticklabels={270,360,540,720,1080},
  log ticks with fixed point, legend pos=outer north east, legend cell align=left]
\addplot[okblue, thick, mark=*] coordinates {(270,\PlotMoonReliefNuOne) (360,\PlotMoonReliefNuTwo) (540,\PlotMoonReliefNuThree) (720,\PlotMoonReliefNuFour) (1080,\PlotMoonReliefNuFive)};
\addlegendentry{Moon, relief}
\addplot[okblue, thick, dashed, mark=o] coordinates {(270,\PlotMoonTextureNuOne) (360,\PlotMoonTextureNuTwo) (540,\PlotMoonTextureNuThree) (720,\PlotMoonTextureNuFour) (1080,\PlotMoonTextureNuFive)};
\addlegendentry{Moon, texture}
\addplot[okvermilion, thick, mark=square*] coordinates {(270,\PlotEarthReliefNuOne) (360,\PlotEarthReliefNuTwo) (540,\PlotEarthReliefNuThree) (720,\PlotEarthReliefNuFour) (1080,\PlotEarthReliefNuFive)};
\addlegendentry{Earth, relief}
\addplot[okvermilion, thick, dashed, mark=square] coordinates {(270,\PlotEarthTextureNuOne) (360,\PlotEarthTextureNuTwo) (540,\PlotEarthTextureNuThree) (720,\PlotEarthTextureNuFour) (1080,\PlotEarthTextureNuFive)};
\addlegendentry{Earth, texture}
\end{loglogaxis}
\end{tikzpicture}
\caption{Structure functions of the primary lunar target window (Site01 of PGDA product 78,
window \texttt{moon-sps-x2-y2}) and the first paired Earth window (Copernicus GLO-30,
\texttt{earth-utm27n-e280-n2880}, Holuhraun), on the \ConfigGridM\,m analysis grid, from the
windows evidence record of the gate G0 pilot. Solid lines are relief (raw heights) and dashed lines
texture (best fit plane removed). Each curve is the equal weight mean over the two grid axes. The
slopes give the Hurst exponents of Table~\ref{tab:first-pair}. The Earth window was chosen because
it contains the decoded radar scene centre, not by similarity.}
\label{fig:structure-functions}
\end{figure}
